El seguimiento del progreso académico de los estudiantes a lo largo de su formación universitaria constituye una herramienta clave para medir la calidad educativa y orientar estrategias de mejora continua. Sin embargo, existe una oportunidad de mejora basada en los datos que cada universidad recolecta de las pruebas Saber 11 y Saber Pro que no es comúnmente aprovechada, y que consiste en tomar toda esta información y transformarla en decisiones de negocio que se reflejen en mejores resultados académicos en los estudiantes.
Esta situación surge dado que la mayoría de las instituciones de educación superior no disponen de sistemas que integren y analicen de manera estructurada la información proveniente de la prueba Saber Pro y se alíe con instituciones escolares que les proporcione información sobre los resultados de las pruebas Saber 11 recientes. A su vez, esta falta de integración impide realizar análisis comparativos que permitan evaluar la evolución de las competencias adquiridas durante el proceso de formación~\cite{garcia2021pruebas}.
En el caso de la Universidad San Buenaventura, donde todavía se cuenta con esta falencia, se muestra como una oportunidad de mejora, contar con un sistema de información que relacione los resultados de ambas pruebas facilitando la identificación de patrones de aprendizaje, fortalezas institucionales y áreas de oportunidad para cada programa académico~\cite{mendieta2024modelo}. La ausencia de este tipo de herramientas limita la capacidad de realizar predicciones sobre el rendimiento futuro de los estudiantes y restringe la generación de estrategias pedagógicas basadas en evidencia~\cite{icfes2024_informeSaber11_2023}.
Por lo tanto, se presenta la necesidad de desarrollar un sistema que consolide los datos de las pruebas Saber Pro, que incorpore técnicas estadísticas y modelos predictivos permitiendo anticipar la fluctuación de las competencias académicas a lo largo del proceso educativo.


